%! TeX root = ../phd-thesis.tex
%! Author = davidedomini

%----------------------------------------------------------------------------------------
\chapter{Deployment on Resource-constrained Edge Devices}
\label{chap:tasksplacement}
\minitoc
%----------------------------------------------------------------------------------------
\section{Motivation}

The learning architectures presented in the previous chapters 
 enable collections of devices to coordinate 
 and learn without relying on a centralized view of the system.
%
Their deployment at the network edge, however, 
 introduces an additional challenge:
 the devices participating in the collective may not have sufficient computational, 
 energy, or memory resources to execute every component of their applications locally.
%
This limitation is becoming increasingly relevant as the number of connected devices 
 and the complexity of edge applications continue to grow.
%
The number of active IoT devices is expected to reach approximately 29.4 billion 
 by 2030\footnote{\url{https://transformainsights.com/news/global-iot-connections-294}}.
%
At this scale, 
 transferring all data and computation to remote cloud infrastructures is often impractical, 
 particularly for applications with strict latency, bandwidth, 
 or privacy requirements~\cite{DBLP:journals/computer/Satyanarayanan17}.

Cloud-edge computing addresses these limitations 
 by distributing computation across a continuum comprising end devices, edge servers, 
 and cloud data centers~\cite{DBLP:journals/jpdc/RosendoCVA22}.
%
Instead of executing an application entirely on a single device or moving it entirely to the cloud,
 its components may be placed at different points of this continuum according to the resources 
 and network conditions currently available.
%
This flexibility gives rise to the \emph{component offloading problem}: 
 determining where each application component should execute in order to balance objectives such as response time, 
 energy consumption, and infrastructure cost.

Offloading decisions become particularly difficult when applications 
 are structured as workflows rather than as indivisible tasks.
%
A workflow may be represented as a directed acyclic graph whose nodes correspond 
 to computational components and whose edges encode data and execution dependencies.
%
The placement of one component therefore affects the cost and feasibility of placing the others:
 executing two dependent components on different devices may reduce local computation 
 but introduce additional communication latency and energy consumption.
%
Component models, architecture description languages, 
 and service-oriented frameworks provide abstractions for decomposing 
 applications~\cite{DBLP:journals/tse/CrnkovicSVC11,DBLP:conf/isola/BallouliBBS18,DBLP:journals/csur/LemosDB16}.
%
However, 
 their deployment is commonly managed through centralized policies, 
  while approaches such as pulverization~\cite{DBLP:journals/fgcs/FarabegoliPCV24} 
  primarily address the decomposition of individual devices and provide limited support 
  for dynamically placing only selected parts of an application.

A further complication is that offloading decisions cannot be optimized independently.
%
When many devices concurrently select the same server or communication path, 
 their individually reasonable choices may collectively produce congestion, 
 resource contention, and server overload.
%
Conversely, 
 distributing requests without accounting for the current state of neighboring devices 
 may leave available resources underutilized.
%
The edge should therefore be understood as a \gls{CAS}~\cite{DBLP:conf/huc/Ferscha15}, 
 in which system-level behavior emerges from the local decisions
 and interactions of multiple heterogeneous entities.
%
An effective offloading policy must account not only for the state of the device making the decision, 
 such as its battery level and computational load, 
 but also for collective information describing the surrounding network and infrastructure.

A large body of work formulates offloading as an optimization problem and addresses it through mathematical programming, 
 heuristic search, genetic algorithms, or machine-learning 
 techniques~\cite{DBLP:journals/comsur/MachB17,DBLP:journals/jnca/LinZCLW20,DBLP:journals/jsa/IslamDGC21,DBLP:journals/cn/SaeikASSDVLAMP21,DBLP:journals/tvt/Al-HabobDAM20}.
%
In particular, 
 reinforcement learning is well suited to dynamic environments because it can learn placement policies directly through interaction, 
 without requiring an exact analytical model of workload and network evolution.
%
Deep reinforcement learning has consequently been applied to offloading in domains including satellite communications, 
 vehicular networks, and crowd management~\cite{DBLP:journals/tvt/ChaiZYXGG23,DBLP:journals/iotj/CuiLZZLQ23,DBLP:journals/tits/WangPLQRWYZZ23}.
%
Nevertheless, 
 existing learning-based approaches commonly reduce offloading to a binary choice between local and remote execution,
  assume a homogeneous infrastructure, or rely on a centralized representation of the global system state.
%
These assumptions limit their applicability to large-scale edge systems, where applications are only partially offloaded, 
 devices and communication links have different roles, and no entity can continuously observe the complete deployment.

Graph-based learning offers a natural foundation for overcoming some of these limitations because 
 both the infrastructure and the application workflow can be represented relationally.
%
However, 
 homogeneous graph representations do not distinguish among end devices, 
 edge servers, cloud servers, communication links, and application dependencies.
%
Moreover, 
 a graph representation alone does not solve the problem of acquiring collective state in a decentralized system.
%
Each agent still requires a compact description of phenomena such as neighborhood density, 
 network congestion, and server utilization, obtained without introducing a central coordinator 
 or assuming global observability.

Accordingly, this chapter addresses the following research question:

\begin{enumerate}[
label=\textbf{RQ\arabic*:},
ref=RQ\arabic*,
leftmargin=*
]
    \item
    How can deep reinforcement learning support decentralized component offloading while accounting 
    for both individual device constraints and collective phenomena such as network congestion and server load?
\label{rq}
\end{enumerate}

To answer this question, 
 the chapter introduces \emph{Informed Deep Heterogeneous Graph Q-Learning} (IDHGQL), 
 a distributed learning approach that combines Deep Q-Learning, \glspl{HGNN}, and aggregate computing.
%
The heterogeneous graph representation preserves the distinct semantics of infrastructure resources, 
 communication links, and application dependencies.
%
Aggregate computing complements this representation by producing decentralized and self-stabilizing summaries 
 of the collective context, allowing each agent to observe relevant system-level phenomena through local interactions.
%
Finally, 
 the macro-component model introduced in previous work~\cite{acsos-model} enables the policy to place individual components 
  of a DAG-structured application rather than choosing only between entirely local and entirely remote execution.
%
The resulting approach uses centralized training and decentralized execution, 
 allowing agents to learn from system-wide experience during training
 while making placement decisions autonomously at runtime.
%
In this way, 
 the chapter extends the thesis investigation from the coordination of distributed learning processes 
 to the resource-aware deployment of their applications across heterogeneous edge infrastructures.

\section{Proposed method: IDHGQL}
\section{Evaluation}
\section{Related Works}
\section{Future Directions}
